\documentclass[11pt]{article}
\usepackage[T1]{fontenc}
\usepackage{lmodern}
\usepackage[margin=1in]{geometry}
\usepackage{amsmath,amssymb,amsthm}
\usepackage{microtype}
\usepackage[hidelinks]{hyperref}
\hypersetup{pdftitle={An Ehrhart--word identity for arbor polytopes},
  pdfauthor={Xiaojian Zeng}}
\newtheorem{theorem}{Theorem}
\newtheorem{lemma}[theorem]{Lemma}
\newtheorem{proposition}[theorem]{Proposition}
\newcommand{\N}{\mathbb Z_{\geq0}}
\newcommand{\R}{\mathbb R}
\newcommand{\Z}{\mathbb Z}
\newcommand{\des}{\operatorname{des}}
\newcommand{\Tle}{T_{\leq}}
\newcommand{\Tge}{T_{\geq}}
\title{An Ehrhart--word identity for arbor polytopes}
\author{Xiaojian Zeng\\[3pt]
  \small PBC School of Finance, Tsinghua University, 100190\\
  \small\texttt{xiaojianzeng@foxmail.com}}
\date{}

\begin{document}
\maketitle
\begin{abstract}
We prove the conjecture of Athanasiadis, Xiao and Yan expressing the
Ehrhart numerator of every arbor polytope as the non-descent enumerator
of words satisfying descendant-set multiplicity conditions. The main step
is a polynomial transformation relating upper bounds on lattice-point
coordinates to lower bounds on binomially weighted word contents at every
subtree. Extracting one coefficient gives an identity for every nonnegative
integer dilation. A decreasing-box construction then yields the word
formula. The argument allows arbitrary nonempty blocks, labels, depth,
and branching.
\end{abstract}

\section{Statement and context}

An \emph{arbor} on $[n]=\{1,\ldots,n\}$ is a rooted tree whose vertices
are the nonempty blocks of a partition of $[n]$ \cite{Chapoton}.
Throughout, $n\geq1$. Write $B_v$ for the block at a vertex $v$, and let
$D_v$ be the union of the blocks in the subtree rooted at $v$, including
$B_v$ itself. Set $b_v=|B_v|$ and $N_v=|D_v|$. The arbor polytope is
\begin{equation}\label{eq:polytope}
 Q_\tau=\left\{x\in\R_{\geq0}^n:
          \sum_{i\in D_v}x_i\leq N_v\text{ for every vertex }v\right\}.
\end{equation}
It contains $[0,1]^n$ and is bounded by the root inequality. It is a
lattice polytope by \cite[Lemma~1.2]{Chapoton}. In particular, its
dimension is $n$. Put $L_\tau(m)=|mQ_\tau\cap\Z^n|$ for $m\in\N$, with
$0Q_\tau=\{0\}$, and define its Ehrhart numerator by
\begin{equation}\label{eq:ehrhart}
 \sum_{m\geq0}L_\tau(m)t^m=
       \frac{h^*(Q_\tau,t)}{(1-t)^{n+1}}.
\end{equation}

Let $W_\tau$ consist of the length-$n$ words $w=w_1\cdots w_n$ over $[n]$
such that
\[
 |\{j:w_j\in D_v\}|\geq N_v\qquad\text{for every vertex }v.
\]
Letters may repeat or be absent. With the usual order on $[n]$, put
\[
 \des(w)=|\{j:1\leq j<n,\ w_j>w_{j+1}\}|,
 \qquad d(w)=n-1-\des(w).
\]
Thus $d(w)$ counts weak rises, including adjacent equal letters.
The set of admissible contents is
\begin{equation}\label{eq:contents}
 A_\tau=\left\{a\in\N^n:\sum_{i=1}^n a_i=n,\quad
         \sum_{i\in D_v}a_i\geq N_v\text{ for every vertex }v\right\}.
\end{equation}

\begin{theorem}\label{thm:main}
For every arbor $\tau$ on $[n]$ and every integer $m\geq0$,
\begin{equation}\label{eq:binomial}
 L_\tau(m)=\sum_{a\in A_\tau}\prod_{i=1}^n\binom{m+1}{a_i},
\end{equation}
where $\binom{m+1}{a_i}=0$ when $a_i>m+1$. Consequently,
\begin{equation}\label{eq:word}
 h^*(Q_\tau,t)=\sum_{w\in W_\tau}t^{n-1-\des(w)}.
\end{equation}
\end{theorem}

Equation~\eqref{eq:word} is Conjecture~5.2 of Athanasiadis, Xiao and
Yan \cite{AXY}. Their Theorem~1.5 proves the case of a root block with
singleton leaves. Athanasiadis and Chapoton \cite[Conjecture~4.4]{AC}
formulate the corresponding assertion for arbitrary preorders; their
Remark~7.3 gives the binomial-content reformulation
\eqref{eq:binomial}. These identities and the fixed-content word identity
used below are therefore prior formulations.

Our contribution is the compatibility of a polynomial transformation with
the two subtree recurrences in Section~\ref{sec:subtrees}. The upper
recurrence is Chapoton's height-counting recurrence
\cite[Proposition~2.1]{Chapoton}. Its relation to a lower-truncation
recurrence proves the asserted identities for every arbor. All algebra
below concerns formal series and finite polynomials.

\section{A truncation identity}

For a series $P$, write $T_{\leq k}P$ for its terms of degrees at most
$k$. For a polynomial $P$, write $T_{\geq k}P$ for its terms of degrees
at least $k$. We first establish the elementary identity that will
intertwine these two operations.

\begin{lemma}\label{lem:local}
For integers $b\geq1$ and $K\geq0$,
\begin{equation}\label{eq:local}
 y^b\sum_{q=0}^{K}\binom{q+b-1}{b-1}(1+y)^{K-q}
       =T_{\geq b}(1+y)^{K+b}.
\end{equation}
\end{lemma}

\begin{proof}
Both sides have degrees between $b$ and $K+b$. For $0\leq j\leq K$,
the coefficient of $y^{b+j}$ on the left is
\begin{align*}
 \sum_{q=0}^{K-j}\binom{q+b-1}{b-1}\binom{K-q}{j}
   &=[z^{K-j}](1-z)^{-b}(1-z)^{-j-1}\\
   &=\binom{K+b}{b+j}.
\end{align*}
The first equality follows by multiplying the two binomial series.
These are exactly the coefficients on the right, including when $K=0$.
\end{proof}

\begin{lemma}\label{lem:operator}
Let $m\geq0$ and $N\geq b\geq1$ be integers. For any polynomial $R(x)$
of degree at most $m(N-b)$, define
\[
 H(x)=T_{\leq mN}\bigl((1-x)^{-b}R(x)\bigr).
\]
Then
\begin{equation}\label{eq:operator}
 \begin{split}
 y^N(1+y)^{mN}H\bigl((1+y)^{-1}\bigr)
   =T_{\geq N}\!\left(
       y^{N-b}(1+y)^{mN+b}R\bigl((1+y)^{-1}\bigr)\right).
 \end{split}
\end{equation}
No sign condition on the coefficients of $R$ is required.
\end{lemma}

\begin{proof}
Expand $R(x)=\sum_s r_s x^s$. For each term put $K=mN-s$; the degree
bound ensures $K\geq mb\geq0$. Its contribution to the left side of
\eqref{eq:operator} is
\begin{align*}
 r_s y^N\sum_{q=0}^{K}\binom{q+b-1}{b-1}(1+y)^{K-q}
   &=r_s y^{N-b}T_{\geq b}(1+y)^{K+b}\\
   &=T_{\geq N}\bigl(r_s y^{N-b}(1+y)^{mN+b-s}\bigr),
\end{align*}
by Lemma~\ref{lem:local}. Summing over $s$ proves the identity.
Every truncated expression on the right is a polynomial. On the left,
$H$ is truncated before substitution and has degree at most $mN$.
Thus the displayed powers clear all denominators; no substitution into
the untruncated infinite series is used.
\end{proof}

\section{Upper and lower subtree counts}\label{sec:subtrees}

Fix $m\in\N$ for this section. Let $C(v)$ be the set of children of $v$.
The structural property of an arbor used in the proof is the disjoint union
\begin{equation}\label{eq:disjoint}
 D_v=B_v\sqcup\bigsqcup_{c\in C(v)}D_c,
 \qquad \sum_{c\in C(v)}N_c=N_v-b_v.
\end{equation}

Define $F_v(x)$ by summing $x^{\sum_{i\in D_v}z_i}$ over the vectors
$z\in\N^{D_v}$ satisfying
\[
 \sum_{i\in D_u}z_i\leq mN_u
 \quad\text{for every descendant }u\text{ of }v,\text{ including }v.
\]
The inequality at $v$ gives $\deg F_v\leq mN_v$. The unrestricted
coordinates in $B_v$ contribute $(1-x)^{-b_v}$. Independence of the child
subtrees in \eqref{eq:disjoint}, followed by the bound at $v$, gives
\begin{equation}\label{eq:F}
 F_v(x)=T_{\leq mN_v}\!\left((1-x)^{-b_v}
                                 \prod_{c\in C(v)}F_c(x)\right).
\end{equation}
This is the upper height recurrence credited above. At the root $r$,
\begin{equation}\label{eq:rootF}
 F_r(1)=L_\tau(m).
\end{equation}
This also holds for $m=0$, since the root bound forces every coordinate
to be zero.

Independently define $G_v(y)$ by summing
\[
 y^{\sum_{i\in D_v}a_i}\prod_{i\in D_v}\binom{m+1}{a_i}
\]
over $a\in\{0,\ldots,m+1\}^{D_v}$ satisfying
\[
 \sum_{i\in D_u}a_i\geq N_u
 \quad\text{for every descendant }u\text{ of }v,\text{ including }v.
\]
The binomial theorem and \eqref{eq:disjoint} give
\begin{equation}\label{eq:G}
 G_v(y)=T_{\geq N_v}\!\left((1+y)^{(m+1)b_v}
                                 \prod_{c\in C(v)}G_c(y)\right).
\end{equation}
The truncation here imposes the lower bound at $v$; all other required
bounds have already been imposed in the child factors.

\begin{proposition}\label{prop:transform}
For every vertex $v$,
\begin{equation}\label{eq:transform}
 G_v(y)=y^{N_v}(1+y)^{mN_v}F_v\bigl((1+y)^{-1}\bigr).
\end{equation}
\end{proposition}

\begin{proof}
Induct from the leaves toward the root. At a vertex $v$, abbreviate
$N=N_v$, $b=b_v$, and $R(x)=\prod_{c\in C(v)}F_c(x)$. The empty product
for a leaf is $1$. In every case,
\[
 \deg R\leq m\sum_{c\in C(v)}N_c=m(N-b).
\]
Using the child identities, or the empty family of identities at a leaf,
we obtain
\begin{align*}
 (1+y)^{(m+1)b}\prod_{c\in C(v)}G_c(y)
   &=y^{N-b}(1+y)^{(m+1)b+m(N-b)}R\bigl((1+y)^{-1}\bigr)\\
   &=y^{N-b}(1+y)^{mN+b}R\bigl((1+y)^{-1}\bigr).
\end{align*}
Apply $T_{\geq N}$ and use Lemma~\ref{lem:operator}. Equations
\eqref{eq:F} and \eqref{eq:G} identify the result as
\eqref{eq:transform}. This proves both the leaf step and the induction.
\end{proof}

At the root, $N_r=n$. The coefficient of $y^n$ in the right side of
\eqref{eq:transform} is the constant term of
$(1+y)^{mn}F_r((1+y)^{-1})$, namely $F_r(1)$. The definition of $G_r$
and \eqref{eq:rootF} therefore imply
\[
 L_\tau(m)=[y^n]G_r(y)
          =\sum_{a\in A_\tau}\prod_{i=1}^n\binom{m+1}{a_i}.
\]
Vectors outside $\{0,\ldots,m+1\}^n$ contribute zero by the binomial
convention. This proves \eqref{eq:binomial} for all $m\geq0$.

\section{From contents to words}

We give the fixed-content conversion explicitly. Its generating-function
form is used in \cite[Section~4]{AXY} and \cite[Remark~7.3]{AC}.

\begin{lemma}\label{lem:words}
Let $a\in\N^n$ with $\sum_i a_i=n$, and let $\mathfrak S_a$ denote the
words containing $a_i$ copies of each letter $i$. Then
\begin{equation}\label{eq:fixedcontent}
 \sum_{m\geq0}\left(\prod_{i=1}^n\binom{m+1}{a_i}\right)t^m
     =\frac{\sum_{w\in\mathfrak S_a}t^{d(w)}}{(1-t)^{n+1}}.
\end{equation}
\end{lemma}

\begin{proof}
For each letter $i$, choose a subset $S_i$ of size $a_i$ of the $m+1$
ordered boxes $0,\ldots,m$. Place $i$ once in each chosen box. Read the
letters in each box in strictly decreasing order, retaining one bar between
consecutive boxes. This produces $n$ letters and exactly $m$ bars.
Empty boxes give initial, final, or adjacent bars. Conversely, split such
a barred word at its bars and record the indices of the boxes containing
each letter. Strict decrease ensures that each letter occurs at most once
per box. The two constructions are inverse.

Fix the unbarred word $w$. A bar is necessary in each gap with
$w_j\leq w_{j+1}$; these are the $d(w)$ required gaps. This condition
is also sufficient, since each resulting box is then strictly decreasing.
After reserving one bar in each required gap, the remaining $m-d(w)$
bars can be placed freely in the $n+1$ gaps, including both ends. Thus
the number of choices above $w$ is zero if $m<d(w)$ and otherwise equals
\[
 \binom{m-d(w)+n}{n}.
\]
Its generating series is
\[
 \sum_{m\geq d(w)}\binom{m-d(w)+n}{n}t^m
       =\frac{t^{d(w)}}{(1-t)^{n+1}}.
\]
Summing over $w\in\mathfrak S_a$ proves \eqref{eq:fixedcontent}.
\end{proof}

\begin{proof}[Completion of the proof of Theorem~\ref{thm:main}]
Sum \eqref{eq:fixedcontent} over the finite set $A_\tau$ and use
\eqref{eq:binomial}. The disjoint union of $\mathfrak S_a$ over
$a\in A_\tau$ is exactly $W_\tau$. Comparing with \eqref{eq:ehrhart}
gives \eqref{eq:word}.
\end{proof}

The proof includes $n=1$, arbitrary leaf-block sizes, and interleaved
labels. At $m=0$ one has $F_v=1$ and $G_v=y^{N_v}$. On the word side
the unique term of degree zero is the decreasing permutation
$n(n-1)\cdots1$, whose content satisfies every lower bound.

\section{Relation to further results and scope}

The associated arbor preorder has the root as its maximum vertex. Each
order ideal is a disjoint union of descendant subtrees. Summing their
inequalities shows that the subtree formulation agrees with the
order-ideal formulation of \cite[Conjecture~4.4]{AC} in this case.
The product decomposition \eqref{eq:disjoint} is essential to the present
proof; the argument does not establish that conjecture for arbitrary
preorders.

For composition polytopes, Athanasiadis \cite[Proposition~4.2]{Athanasiadis}
gives an interpretation using modified ascents, with an initial letter
$w_0=1$. The corresponding minimum-vertex extension appears in
\cite[Proposition~4.5]{AC}. Those statistics differ from the ordinary
non-descents in Theorem~\ref{thm:main}. Yan
\cite[Corollary~5.4]{Yan} gives a word formula for the related reflexive
polytope $Q_\tau+\Delta_E-\mathbf1$, where
$\Delta_E=\operatorname{conv}(0,e_i:i\in E)$; that formula uses an
auxiliary letter. It concerns a different polytope. These comparisons
identify the precise assertion proved here without a claim of exhaustive
priority. No equivalent proof was found in the primary passages inspected
through 14 September 2026.

Theorem~\ref{thm:main} also makes no assertion about real-rootedness or
about the support enumerator
$\sum_{x\in Q_\tau\cap\Z^n}t^{|\{i:x_i>0\}|}$, which counts points
in a single undilated polytope. For general preorder polytopes, the
real-rootedness of this support enumerator was proposed in
\cite[Conjecture~5.3]{AC}. Zijian Zeng
\cite[Proposition~3.1 and Theorem~4.2]{ZengSupport} certifies an
eight-element height-two poset counterexample using exact support
enumeration and a rational non-real-root certificate. We use the author
manuscript dated 31 August 2026 for this comparison. That result concerns
the support enumerator, and is not a counterexample to
Theorem~\ref{thm:main}.

\paragraph{Exact checks and computational assistance.}
The proof is independent of computation. Two separately implemented
Python programs using exact integer arithmetic checked the truncation
identities, the box construction, and direct coordinate and word counts.
The direct enumeration includes all rooted trees on the blocks of every
partition of $[n]$ for $1\leq n\leq4$, with $0\leq m\leq n$:
153 arbors and 740 arbor--dilation pairs. Three specified six-label arbors
and one seven-label arbor were also checked for $m=0,1,2,3$.
These finite checks test conventions and formulas; they are not an
extension of the previously reported size-seven preorder tests
\cite[Section~4.3]{AC}, and do not replace the proof.
OpenAI language-model assistance was used in developing the argument,
writing the verification programs, checking primary-source comparisons,
and preparing this manuscript. A separate model context reviewed the
argument and supplied a separate enumeration implementation. This is not
formal proof-assistant certification or journal peer review.

\bibliographystyle{amsplain}

\raggedright
\bibliography{references}

\end{document}
